\documentclass[UTF8,zihao=-4]{ctexart}
\usepackage{amsmath,amssymb,mathtools}
\usepackage[a4paper,margin=2.2cm]{geometry}
\setlength{\parindent}{0pt}
\setlength{\parskip}{0.45em}
\allowdisplaybreaks

\begin{document}

\section*{一}

两端固定、长度为 $l$ 的琴弦，忽略重力。初始时刻在琴弦的 $x=c$ 处，拉离平衡位置 $h$ 的高度，然后静止释放。求该体系的定解。

定解问题为
\[
u_{tt}-a^2u_{xx}=0,\qquad u(0,t)=u(l,t)=0,
\]
\[
u(x,0)=f(x)=
\begin{cases}
\dfrac{h}{c}x,&0\le x\le c,\\[4pt]
\dfrac{h}{l-c}(l-x),&c\le x\le l,
\end{cases}
\qquad u_t(x,0)=0.
\]

由边界条件可知
\[
X_n(x)=\sin\dfrac{n\pi x}{l},\qquad \omega_n=\dfrac{an\pi}{l},
\qquad n=1,2,\dots
\]
\[
u(x,t)=\sum_{n=1}^{\infty}(A_n\cos\omega_n t+B_n\sin\omega_n t)
\sin\dfrac{n\pi x}{l}.
\]
由 $u_t(x,0)=0$，得 $B_n=0$。又
\[
A_n=\dfrac{2}{l}\left[
\int_0^c\dfrac{h}{c}x\sin\dfrac{n\pi x}{l}\,\mathrm{d}x
+\int_c^l\dfrac{h}{l-c}(l-x)\sin\dfrac{n\pi x}{l}\,\mathrm{d}x
\right]
=\dfrac{2hl^2}{n^2\pi^2c(l-c)}\sin\dfrac{n\pi c}{l}.
\]
故
\[
u(x,t)=\sum_{n=1}^{\infty}
\dfrac{2hl^2}{n^2\pi^2c(l-c)}\sin\dfrac{n\pi c}{l}
\cos\dfrac{an\pi t}{l}\sin\dfrac{n\pi x}{l}.
\]

\section*{二}

在一维系统 $(0\le x\le l)$ 上，求解定解问题：
\[
u_{tt}-a^2u_{xx}=0,\qquad
u(0,t)=0,\quad u_x(l,t)=0,
\]
\[
u(x,0)=cx,\qquad u_t(x,0)=0.
\]

由边界条件可知
\[
k_n=\dfrac{(2n+1)\pi}{2l},\qquad X_n(x)=\sin k_nx,\qquad
\omega_n=ak_n,\qquad n=0,1,2,\dots
\]
\[
u(x,t)=\sum_{n=0}^{\infty}(A_n\cos\omega_nt+B_n\sin\omega_nt)\sin k_nx.
\]
由 $u_t(x,0)=0$，得 $B_n=0$。又
\[
A_n=\dfrac{2}{l}\int_0^lcx\sin k_nx\,\mathrm{d}x
=\dfrac{8cl(-1)^n}{(2n+1)^2\pi^2}.
\]
故
\[
u(x,t)=\sum_{n=0}^{\infty}
\dfrac{8cl(-1)^n}{(2n+1)^2\pi^2}
\cos\dfrac{(2n+1)\pi at}{2l}\sin\dfrac{(2n+1)\pi x}{2l}.
\]

\section*{三}

利用分离变数法，在一维系统 $(0\le x\le l)$ 上，求下列定解问题：
\[
u_{tt}+u_t-a^2u_{xx}=0,\qquad
u_x(0,t)=u_x(l,t)=0,
\]
\[
u(x,0)=\sin x,\qquad u_t(x,0)=0.
\]

由边界条件可知
\[
k_n=\dfrac{n\pi}{l},\qquad X_n(x)=\cos k_nx,\qquad
\omega_n=ak_n,\qquad n=0,1,2,\dots
\]
令
\[
u(x,t)=\sum_{n=0}^{\infty}c_n\Phi_n(t)\cos k_nx,\qquad
\Phi_n(0)=1,\quad \Phi_n'(0)=0.
\]
其中
\[
c_0=\dfrac{1}{l}\int_0^l\sin x\,\mathrm{d}x=\dfrac{1-\cos l}{l},
\]
\[
c_n=\dfrac{2}{l}\int_0^l\sin x\cos\dfrac{n\pi x}{l}\,\mathrm{d}x
=\dfrac{2l\left[1-(-1)^n\cos l\right]}{l^2-n^2\pi^2},\qquad n\ge1.
\]
$\Phi_n$ 满足
\[
\Phi_n''+\Phi_n'+\omega_n^2\Phi_n=0.
\]
记
\[
r_{n\pm}=\dfrac{-1\pm\sqrt{1-4\omega_n^2}}{2}.
\]
则
\[
\Phi_n(t)=
\begin{cases}
\dfrac{-r_{n-}e^{r_{n+}t}+r_{n+}e^{r_{n-}t}}{r_{n+}-r_{n-}},
&\omega_n\ne\dfrac12,\\[8pt]
\left(1+\dfrac{t}{2}\right)e^{-t/2},&\omega_n=\dfrac12.
\end{cases}
\]
故
\[
u(x,t)=c_0\Phi_0(t)+\sum_{n=1}^{\infty}
c_n\Phi_n(t)\cos\dfrac{n\pi x}{l}.
\]

\section*{四}

在一维系统 $(0\le x\le l)$ 上，求下列定解问题：
\[
u_t-a^2u_{xx}=0,\qquad
u_x(0,t)=0,\quad u(l,t)=0,
\]
\[
u(x,0)=x.
\]

由边界条件可知
\[
k_n=\dfrac{(2n+1)\pi}{2l},\qquad X_n(x)=\cos k_nx,\qquad n=0,1,2,\dots
\]
\[
u(x,t)=\sum_{n=0}^{\infty}A_ne^{-a^2k_n^2t}\cos k_nx.
\]
由 $u(x,0)=x$，
\[
A_n=\dfrac{2}{l}\int_0^lx\cos k_nx\,\mathrm{d}x
=\dfrac{4l(-1)^n}{(2n+1)\pi}
-\dfrac{8l}{(2n+1)^2\pi^2}.
\]
故
\[
u(x,t)=\sum_{n=0}^{\infty}
\left[
\dfrac{4l(-1)^n}{(2n+1)\pi}
-\dfrac{8l}{(2n+1)^2\pi^2}
\right]
e^{-a^2(2n+1)^2\pi^2t/(4l^2)}
\cos\dfrac{(2n+1)\pi x}{2l}.
\]

\section*{五}

在方形区域 $(0\le x\le a,\ 0\le y\le b)$ 中，求下列定解问题：
\[
u_{xx}+u_{yy}=0,
\]
\[
u(0,y)=0,\quad u(a,y)=c(1-y),\qquad
u(x,0)=c,\quad u(x,b)=0.
\]

令 $u=v+w$，其中
\[
v(0,y)=v(x,0)=v(x,b)=0,\qquad v(a,y)=c(1-y),
\]
\[
w(0,y)=w(a,y)=w(x,b)=0,\qquad w(x,0)=c.
\]
则
\[
v=\sum_{n=1}^{\infty}A_n\sinh\dfrac{n\pi x}{b}\sin\dfrac{n\pi y}{b},
\]
\[
A_n\sinh\dfrac{n\pi a}{b}
=\dfrac{2}{b}\int_0^bc(1-y)\sin\dfrac{n\pi y}{b}\,\mathrm{d}y
=\dfrac{2c\left[1+(b-1)(-1)^n\right]}{n\pi}.
\]
又
\[
w=\sum_{m=1}^{\infty}B_m\sin\dfrac{m\pi x}{a}\sinh\dfrac{m\pi(b-y)}{a},
\]
\[
B_m\sinh\dfrac{m\pi b}{a}
=\dfrac{2}{a}\int_0^ac\sin\dfrac{m\pi x}{a}\,\mathrm{d}x
=\dfrac{2c\left[1-(-1)^m\right]}{m\pi}.
\]
故
\[
u=\sum_{n=1}^{\infty}
\dfrac{2c\left[1+(b-1)(-1)^n\right]}{n\pi\sinh\dfrac{n\pi a}{b}}
\sinh\dfrac{n\pi x}{b}\sin\dfrac{n\pi y}{b}
\]
\[
\sum_{m=1}^{\infty}
\dfrac{2c\left[1-(-1)^m\right]}{m\pi\sinh\dfrac{m\pi b}{a}}
\sin\dfrac{m\pi x}{a}\sinh\dfrac{m\pi(b-y)}{a}.
\]

\section*{六}

在方形区域 $(0\le x\le a,\ 0\le y\le b)$ 中，求下列定解问题：
\[
u_{xx}+u_{yy}=0,
\]
\[
u_x(0,y)=0,\quad u(a,y)=2y,\qquad
u_y(x,0)=u_y(x,b)=0.
\]

由 $y$ 方向边界条件可知
\[
Y_n(y)=\cos\dfrac{n\pi y}{b},\qquad n=0,1,2,\dots
\]
故
\[
u=A_0+\sum_{n=1}^{\infty}A_n\cosh\dfrac{n\pi x}{b}\cos\dfrac{n\pi y}{b}.
\]
由 $u(a,y)=2y$，
\[
A_0=\dfrac1b\int_0^b2y\,\mathrm{d}y=b,
\]
\[
A_n\cosh\dfrac{n\pi a}{b}
=\dfrac{2}{b}\int_0^b2y\cos\dfrac{n\pi y}{b}\,\mathrm{d}y
=\dfrac{4b\left[(-1)^n-1\right]}{n^2\pi^2}.
\]
故
\[
u(x,y)=b+\sum_{n=1}^{\infty}
\dfrac{4b\left[(-1)^n-1\right]}{n^2\pi^2\cosh\dfrac{n\pi a}{b}}
\cosh\dfrac{n\pi x}{b}\cos\dfrac{n\pi y}{b}.
\]

\section*{七}

在一个内半径为 $\rho_1$，外半径为 $\rho_2$ 的圆环
$(\rho_1\le\rho\le\rho_2,\ 0\le\theta<2\pi)$ 内，求解拉普拉斯方程的定解：
\[
\dfrac1\rho\dfrac{\partial}{\partial\rho}\left(\rho\dfrac{\partial u}{\partial\rho}\right)
+\dfrac1{\rho^2}\dfrac{\partial^2u}{\partial\theta^2}=0,
\]
\[
u(\rho_1,\theta)=\sin\theta,\qquad u_\rho(\rho_2,\theta)=0.
\]

只含 $\sin\theta$ 项，设
\[
u(\rho,\theta)=\left(C\rho+\dfrac{D}{\rho}\right)\sin\theta.
\]
由边界条件
\[
C\rho_1+\dfrac{D}{\rho_1}=1,\qquad
C-\dfrac{D}{\rho_2^2}=0.
\]
解得
\[
C=\dfrac{\rho_1}{\rho_1^2+\rho_2^2},\qquad
D=\dfrac{\rho_1\rho_2^2}{\rho_1^2+\rho_2^2}.
\]
故
\[
u(\rho,\theta)=
\dfrac{\rho_1}{\rho_1^2+\rho_2^2}
\left(\rho+\dfrac{\rho_2^2}{\rho}\right)\sin\theta.
\]

\section*{八}

在一个内半径为 $\rho_1$，外半径为 $\rho_2$ 的圆环
$(\rho_1\le\rho\le\rho_2,\ 0\le\theta<2\pi)$ 内，求解拉普拉斯方程的定解：
\[
\dfrac1\rho\dfrac{\partial}{\partial\rho}\left(\rho\dfrac{\partial u}{\partial\rho}\right)
+\dfrac1{\rho^2}\dfrac{\partial^2u}{\partial\theta^2}
=12\rho^2\cos^2\theta,
\]
\[
u(\rho_1,\theta)=0,\qquad u_\rho(\rho_2,\theta)=0.
\]

因为
\[
12\rho^2\cos^2\theta=6\rho^2+6\rho^2\cos2\theta,
\]
设
\[
u(\rho,\theta)=R_0(\rho)+R_2(\rho)\cos2\theta.
\]
则
\[
R_0''+\dfrac1\rho R_0'=6\rho^2,\qquad
R_0(\rho_1)=0,\quad R_0'(\rho_2)=0,
\]
\[
R_0(\rho)=\dfrac38(\rho^4-\rho_1^4)-\dfrac32\rho_2^4\ln\dfrac{\rho}{\rho_1}.
\]
又
\[
R_2''+\dfrac1\rho R_2'-\dfrac4{\rho^2}R_2=6\rho^2,
\qquad R_2(\rho_1)=0,\quad R_2'(\rho_2)=0.
\]
设
\[
R_2(\rho)=\dfrac12\rho^4+C\rho^2+\dfrac{D}{\rho^2}.
\]
由边界条件得
\[
C=-\dfrac{2\rho_2^6+\rho_1^6}{2(\rho_1^4+\rho_2^4)},\qquad
D=\dfrac{\rho_1^4\rho_2^4(2\rho_2^2-\rho_1^2)}
{2(\rho_1^4+\rho_2^4)}.
\]
故
\[
u(\rho,\theta)=
\dfrac38(\rho^4-\rho_1^4)-\dfrac32\rho_2^4\ln\dfrac{\rho}{\rho_1}
+\left(\dfrac12\rho^4+C\rho^2+\dfrac{D}{\rho^2}\right)\cos2\theta.
\]

\section*{九}

在半径为 $a$ 的半圆盘中 $(0\le\rho\le a,\ 0\le\theta\le\pi)$，求解定解问题：
\[
\dfrac1\rho\dfrac{\partial}{\partial\rho}\left(\rho\dfrac{\partial u}{\partial\rho}\right)
+\dfrac1{\rho^2}\dfrac{\partial^2u}{\partial\theta^2}=0,
\]
\[
u(\rho,0)=0,\qquad u(\rho,\pi)=0,\qquad
u(a,\theta)=u_0\theta(\pi-\theta).
\]

设
\[
u(\rho,\theta)=\sum_{n=1}^{\infty}A_n\rho^n\sin n\theta.
\]
由 $\rho=a$，
\[
u_0\theta(\pi-\theta)=\sum_{n=1}^{\infty}A_na^n\sin n\theta.
\]
因此
\[
A_na^n=\dfrac{2u_0}{\pi}\int_0^\pi \theta(\pi-\theta)\sin n\theta\,\mathrm{d}\theta
=\dfrac{4u_0\left[1-(-1)^n\right]}{\pi n^3}.
\]
故
\[
u(\rho,\theta)=
\sum_{n=1}^{\infty}
\dfrac{4u_0\left[1-(-1)^n\right]}{\pi n^3}
\left(\dfrac{\rho}{a}\right)^n\sin n\theta.
\]

\section*{十}

在一维系统 $(0\le x\le l)$ 上，求解定解问题：
\[
u_{tt}-a^2u_{xx}=b\sinh x,\qquad
u(0,t)=u(l,t)=0,
\]
\[
u(x,0)=0,\qquad u_t(x,0)=0.
\]

令
\[
u(x,t)=\sum_{n=1}^{\infty}T_n(t)\sin\dfrac{n\pi x}{l},\qquad
\omega_n=\dfrac{an\pi}{l}.
\]
右端展开为
\[
b\sinh x=\sum_{n=1}^{\infty}F_n\sin\dfrac{n\pi x}{l},
\]
\[
F_n=\dfrac{2b}{l}\int_0^l\sinh x\sin\dfrac{n\pi x}{l}\,\mathrm{d}x
=\dfrac{2b(-1)^{n+1}n\pi\sinh l}{l^2+n^2\pi^2}.
\]
故
\[
T_n''+\omega_n^2T_n=F_n,\qquad T_n(0)=T_n'(0)=0,
\]
\[
T_n(t)=\dfrac{F_n}{\omega_n^2}(1-\cos\omega_nt).
\]
即
\[
u(x,t)=\sum_{n=1}^{\infty}
\dfrac{2bl^2(-1)^{n+1}\sinh l}{a^2n\pi(l^2+n^2\pi^2)}
\left(1-\cos\dfrac{an\pi t}{l}\right)\sin\dfrac{n\pi x}{l}.
\]

\section*{十一}

在方形区域 $(0\le x\le a,\ 0\le y\le b)$ 中，求下列定解问题：
\[
u_{xx}+u_{yy}=-1,
\]
\[
u(0,y)=2,\quad u(a,y)=2,\qquad
u_y(x,0)=2,\quad u_y(x,b)=2.
\]

令 $u=v+w$，列表为
\[
\begin{array}{c|cc|cc|c}
 & x=0 & x=a & y=0 & y=b & \text{方程}\\ \hline
v & 2 & 2 & v_y=0 & v_y=0 & v_{xx}+v_{yy}=-1\\
w & 0 & 0 & w_y=2 & w_y=2 & w_{xx}+w_{yy}=0
\end{array}
\]
先解 $v$，取
\[
v(x,y)=2+\dfrac{a}{2}x-\dfrac{x^2}{2}.
\]
再解 $w$：
\[
w_{xx}+w_{yy}=0,\qquad w(0,y)=w(a,y)=0,\qquad
w_y(x,0)=w_y(x,b)=2.
\]
设
\[
w=\sum_{n=1}^{\infty}Y_n(y)\sin\dfrac{n\pi x}{a},\qquad
\mu_n=\dfrac{n\pi}{a}.
\]
常数 $2$ 的正弦展开系数为
\[
d_n=\dfrac{2}{a}\int_0^a2\sin\dfrac{n\pi x}{a}\,\mathrm{d}x
=\dfrac{4[1-(-1)^n]}{n\pi}.
\]
取
\[
Y_n(y)=\dfrac{d_n}{\mu_n\cosh\dfrac{\mu_nb}{2}}
\sinh\mu_n\left(y-\dfrac b2\right).
\]
故
\[
u(x,y)=2+\dfrac{a}{2}x-\dfrac{x^2}{2}
+\sum_{n=1}^{\infty}
\dfrac{4a[1-(-1)^n]}{n^2\pi^2\cosh\dfrac{n\pi b}{2a}}
\sinh\dfrac{n\pi}{a}\left(y-\dfrac b2\right)
\sin\dfrac{n\pi x}{a}.
\]

\section*{十二}

在一维系统 $(0\le x\le l)$ 上，求解定解问题 $(g\text{ 是常数})$：
\[
u_{tt}-a^2u_{xx}=g,\qquad
u(0,t)=0,\quad u_x(l,t)=0,
\]
\[
u(x,0)=0,\qquad u_t(x,0)=0.
\]

由边界条件可知
\[
k_n=\dfrac{(2n+1)\pi}{2l},\qquad X_n(x)=\sin k_nx,\qquad
\omega_n=ak_n.
\]
令
\[
u=\sum_{n=0}^{\infty}T_n(t)\sin k_nx.
\]
右端展开为
\[
g=\sum_{n=0}^{\infty}G_n\sin k_nx,\qquad
G_n=\dfrac{2}{l}\int_0^lg\sin k_nx\,\mathrm{d}x
=\dfrac{2g}{lk_n}.
\]
于是
\[
T_n''+\omega_n^2T_n=G_n,\qquad T_n(0)=T_n'(0)=0,
\]
\[
T_n(t)=\dfrac{G_n}{\omega_n^2}(1-\cos\omega_nt).
\]
故
\[
u(x,t)=\sum_{n=0}^{\infty}
\dfrac{16gl^2}{a^2(2n+1)^3\pi^3}
\left(1-\cos\dfrac{(2n+1)\pi at}{2l}\right)
\sin\dfrac{(2n+1)\pi x}{2l}.
\]

\section*{十三}

在一维系统 $(0\le x\le l)$ 上，求解定解问题：
\[
u_{tt}-a^2u_{xx}=A\sin\omega t,\qquad
u(0,t)=0,\quad u_x(l,t)=0,
\]
\[
u(x,0)=0,\qquad u_t(x,0)=0.
\]

令
\[
k_n=\dfrac{(2n+1)\pi}{2l},\qquad \Omega_n=ak_n,\qquad
u=\sum_{n=0}^{\infty}T_n(t)\sin k_nx.
\]
于是
\[
T_n''+\Omega_n^2T_n=\dfrac{2A}{lk_n}\sin\omega t,\qquad
T_n(0)=T_n'(0)=0.
\]
当 $\omega\ne\Omega_n$ 时，
\[
u(x,t)=\sum_{n=0}^{\infty}
\dfrac{2A}{lk_n(\Omega_n^2-\omega^2)}
\left(\sin\omega t-\dfrac{\omega}{\Omega_n}\sin\Omega_nt\right)\sin k_nx.
\]
当 $\omega=\Omega_m$ 时，第 $m$ 项改为
\[
\dfrac{A}{lk_m\Omega_m^2}
\left(\sin\Omega_mt-\Omega_mt\cos\Omega_mt\right)\sin k_mx.
\]

\section*{十四}

在一维系统 $(0\le x\le l)$ 上，求解定解问题：
\[
u_t-a^2u_{xx}=A\sin\omega t,\qquad
u(0,t)=0,\quad u_x(l,t)=0,
\]
\[
u(x,0)=A\cos\dfrac{\pi x}{l}.
\]

令
\[
k_n=\dfrac{(2n+1)\pi}{2l},\qquad
\lambda_n=a^2k_n^2,\qquad
u=\sum_{n=0}^{\infty}T_n(t)\sin k_nx.
\]
初值系数为
\[
C_n=\dfrac{2A}{l}\int_0^l\cos\dfrac{\pi x}{l}\sin k_nx\,\mathrm{d}x
=\dfrac{4A(2n+1)}{\pi(2n-1)(2n+3)}.
\]
于是
\[
T_n'+\lambda_nT_n=\dfrac{2A}{lk_n}\sin\omega t,\qquad T_n(0)=C_n.
\]
故
\[
\begin{aligned}
u(x,t)=\sum_{n=0}^{\infty}\Bigg[
&\dfrac{4A(2n+1)}{\pi(2n-1)(2n+3)}e^{-\lambda_nt}\\
&+\dfrac{2A}{lk_n(\lambda_n^2+\omega^2)}
\left(\lambda_n\sin\omega t-\omega\cos\omega t+\omega e^{-\lambda_nt}\right)
\Bigg]\sin k_nx.
\end{aligned}
\]

\section*{十五}

在一维系统 $(0\le x\le l)$ 上，求解定解问题：
\[
u_{tt}-a^2u_{xx}=bx(l-x),
\]
\[
u(0,t)=0,\quad u(l,t)=A\sin\omega t,\qquad
u(x,0)=0,\quad u_t(x,0)=0.
\]

令
\[
u(x,t)=\dfrac{Ax}{l}\sin\omega t+v(x,t),
\qquad
v=\sum_{n=1}^{\infty}T_n(t)\sin\dfrac{n\pi x}{l}.
\]
记
\[
k_n=\dfrac{n\pi}{l},\qquad \Omega_n=ak_n.
\]
则
\[
T_n''+\Omega_n^2T_n=
\dfrac{4bl^2[1-(-1)^n]}{n^3\pi^3}
-\dfrac{2A\omega^2(-1)^n}{n\pi}\sin\omega t,
\]
\[
T_n(0)=0,\qquad
T_n'(0)=\dfrac{2A\omega(-1)^n}{n\pi}.
\]
当 $\omega\ne\Omega_n$ 时，
\[
\begin{aligned}
u(x,t)=\dfrac{Ax}{l}\sin\omega t+\sum_{n=1}^{\infty}\Bigg[
&\dfrac{4bl^4[1-(-1)^n]}{a^2n^5\pi^5}(1-\cos\Omega_nt)\\
&+\dfrac{2A\omega l(-1)^n}{an^2\pi^2}\sin\Omega_nt\\
&-\dfrac{2A\omega^2(-1)^n}{n\pi(\Omega_n^2-\omega^2)}
\left(\sin\omega t-\dfrac{\omega}{\Omega_n}\sin\Omega_nt\right)
\Bigg]\sin k_nx.
\end{aligned}
\]
当 $\omega=\Omega_m$ 时，第 $m$ 项中最后一项改为
\[
-\dfrac{A(-1)^m}{m\pi}
\left(\sin\Omega_mt-\Omega_mt\cos\Omega_mt\right).
\]

\section*{十六}

在一维系统 $(0\le x\le l)$ 上，求下列定解问题：
\[
u_t-a^2u_{xx}=0,\qquad
u_x(0,t)=0,\quad u_x(l,t)=Be^{-\alpha t},
\]
\[
u(x,0)=0.
\]

记
\[
\mu=\dfrac{\sqrt{\alpha}}{a}.
\]
取
\[
p(x,t)=-\dfrac{B}{\mu\sin\mu l}\cos\mu x\,e^{-\alpha t}.
\]
令 $u=p+w$，则
\[
w_t-a^2w_{xx}=0,\qquad w_x(0,t)=w_x(l,t)=0,
\]
\[
w(x,0)=\dfrac{B}{\mu\sin\mu l}\cos\mu x.
\]
设
\[
w=C_0+\sum_{n=1}^{\infty}C_ne^{-a^2n^2\pi^2t/l^2}\cos\dfrac{n\pi x}{l}.
\]
则
\[
C_0=\dfrac{1}{l}\int_0^l\dfrac{B}{\mu\sin\mu l}\cos\mu x\,\mathrm{d}x
=\dfrac{B}{\mu^2l},
\]
\[
C_n=\dfrac{2}{l}\int_0^l\dfrac{B}{\mu\sin\mu l}
\cos\mu x\cos\dfrac{n\pi x}{l}\,\mathrm{d}x
=\dfrac{2B(-1)^n}{l\left(\mu^2-\dfrac{n^2\pi^2}{l^2}\right)}.
\]
故
\[
u(x,t)=-\dfrac{B}{\mu\sin\mu l}\cos\mu x\,e^{-\alpha t}
+\dfrac{B}{\mu^2l}
+\sum_{n=1}^{\infty}
\dfrac{2B(-1)^n}{l\left(\mu^2-\dfrac{n^2\pi^2}{l^2}\right)}
e^{-a^2n^2\pi^2t/l^2}\cos\dfrac{n\pi x}{l}.
\]

\section*{十七}

在一维系统 $(0\le x\le l)$ 上，求下列定解问题：
\[
u_{tt}-a^2u_{xx}=0,
\]
\[
u(0,t)=\cos\dfrac{\pi at}{l},\qquad u_x(l,t)=0,
\]
\[
u(x,0)=\cos\dfrac{\pi x}{l},\qquad
u_t(x,0)=\sin\dfrac{\pi x}{2l}.
\]

取
\[
p(x,t)=\cos\dfrac{\pi x}{l}\cos\dfrac{\pi at}{l},
\qquad u=p+v.
\]
则
\[
v_{tt}-a^2v_{xx}=0,\qquad v(0,t)=0,\quad v_x(l,t)=0,
\]
\[
v(x,0)=0,\qquad v_t(x,0)=\sin\dfrac{\pi x}{2l}.
\]
由边界条件
\[
k_n=\dfrac{(2n+1)\pi}{2l},\qquad \omega_n=ak_n.
\]
因为初速只含 $n=0$ 项，
\[
v(x,t)=\dfrac{1}{\omega_0}\sin\omega_0t\sin k_0x
=\dfrac{2l}{a\pi}\sin\dfrac{a\pi t}{2l}\sin\dfrac{\pi x}{2l}.
\]
故
\[
u(x,t)=\cos\dfrac{\pi x}{l}\cos\dfrac{\pi at}{l}
+\dfrac{2l}{a\pi}\sin\dfrac{a\pi t}{2l}\sin\dfrac{\pi x}{2l}.
\]

\section*{十八}

在一维系统 $(0\le x\le l)$ 上，求下列定解问题：
\[
u_t-a^2u_{xx}=0,
\]
\[
u(0,t)=Ae^{-\alpha t},\qquad u_x(l,t)=Be^{-\beta t},
\]
\[
u(x,0)=0.
\]

记
\[
\mu=\dfrac{\sqrt{\alpha}}{a},\qquad
\nu=\dfrac{\sqrt{\beta}}{a}.
\]
取
\[
p(x,t)=A\dfrac{\cos\mu(l-x)}{\cos\mu l}e^{-\alpha t}
+\dfrac{B}{\nu\cos\nu l}\sin\nu x\,e^{-\beta t}.
\]
令 $u=p+v$，则
\[
v_t-a^2v_{xx}=0,\qquad v(0,t)=0,\quad v_x(l,t)=0,
\]
\[
v(x,0)=-
A\dfrac{\cos\mu(l-x)}{\cos\mu l}
-\dfrac{B}{\nu\cos\nu l}\sin\nu x.
\]
设
\[
k_n=\dfrac{(2n+1)\pi}{2l},\qquad \lambda_n=a^2k_n^2,
\qquad
v=\sum_{n=0}^{\infty}C_ne^{-\lambda_nt}\sin k_nx.
\]
则
\[
C_n=-\dfrac{2}{l}
\left[
\dfrac{Ak_n}{k_n^2-\mu^2}
+\dfrac{B}{k_n^2-\nu^2}
\right].
\]
故
\[
u(x,t)=
A\dfrac{\cos\mu(l-x)}{\cos\mu l}e^{-\alpha t}
+\dfrac{B}{\nu\cos\nu l}\sin\nu x\,e^{-\beta t}
+\sum_{n=0}^{\infty}C_ne^{-\lambda_nt}\sin k_nx.
\]

\end{document}
